\subsection{Bundles of alternating multilinear maps}

In nos. 7.8.1 to 7.8.5, we assume that K has characteristic 0 or that the vector bundles under consideration have finite rank. It is not known whether the definition of the functor \(\alpha_{d}\) given in 7.8.1 can be done without any restriction.

7.8.1. Let us take \(\mathbf{I}_{+} = \{0\}\) , \(\mathbf{I}_{-} = \{1\}\) and let \(d\) be an integer \(\geqslant 1\) . We define a

\(^{1}\) When \(M\) is of finite rank, we write \(M^{*}\) in place of \(M'\).

vector functor \(\alpha_{d}\) by denoting by \(\alpha_{d}(\mathcal{V})\) the Banach space of continuous alternating \(d\)-linear maps from \(\mathbf{V}_{1}^{d}\) into \(\mathbf{V}_{0}\) and setting \(\alpha_{d}(\mathbf{f})(u) = f_{0} \circ u \circ f_{1}^{d}\) for \(u \in \alpha_{d}(\mathcal{V})\) . The vector bundle \(\alpha_{d}((\mathbf{M}_{1}, \mathbf{M}_{0}))\) is denoted \(\operatorname{Alt}^{d}(\mathbf{M}_{1}; \mathbf{M}_{0})\) and is called the vector bundle of alternating \(d\)-linear maps from \(\mathbf{M}_{1}\) to \(\mathbf{M}_{0}\) .

The canonical injection of \(\operatorname{Alt}^{d}(\mathbf{M}_{1};\mathbf{M}_{0})\) to \(\mathcal{L}(\mathbf{M}_{1},\ldots,\mathbf{M}_{1};\mathbf{M}_{0})\) is a morphism of vector bundles; \(\operatorname{Alt}^{d}(\mathbf{M}_{1};\mathbf{M}_{0})\) is a vector subbundle of \(\mathcal{L}(\mathbf{M}_{1},\ldots,\mathbf{M}_{1};\mathbf{M}_{0})\) .

One has \(\operatorname{Alt}^1 (\mathbf{M}_1;\mathbf{M}_0) = \mathcal{L}(\mathbf{M}_1;\mathbf{M}_0)\) . We set \(\operatorname{Alt}^0 (\mathbf{M}_1;\mathbf{M}_0) = \mathbf{M}_0\) .

If \(\omega\) is a section \(^{1}\) of \(\mathsf{Alt}^{d}(\mathbf{M}_{1};\mathbf{M}_{0})\) and if \(s_{1},\ldots,s_{d}\) are sections of \(\mathbf{M}_{1}\) there exists one and only one section of \(\mathbf{M}_{0}\) , denoted \(\omega(s_{1},\ldots,s_{d})\) , such that

\[
\omega (s _ {1}, \dots , s _ {d}) (b) = \omega (b) (s _ {1} (b), \dots , s _ {d} (b)) \quad \text { for all } b \in \mathbf {B}.
\]

7.8.2. Let \(\varphi\) be a pairing of \(N \times_{B} N'\) to \(N''\) (cf. no. 7.3.1); for each \(b\) , we have a bilinear map \(\varphi_b\) of \(N_b \times N_b'\) to \(N_b''\) which defines (cf. A, III, p. 142) a bilinear map of

\[
\operatorname{Alt} ^ {d} (\mathbf {M}; \mathbf {N}) _ {b} \times \operatorname{Alt} ^ {e} (\mathbf {M}; \mathbf {N} ^ {\prime}) _ {b}
\]

to \(\operatorname{Alt}^{d + e}(\mathbf{M};\mathbf{N}^{\prime \prime})_{b}\) . The collection of these bilinear maps defines a pairing \(u\) of

\[
\operatorname{Alt} ^ {d} (\mathbf {M}; \mathbf {N}) \times_ {\mathbf {B}} \operatorname{Alt} ^ {e} (\mathbf {M}; \mathbf {N} ^ {\prime})
\]

to \(\mathsf{Alt}^{d + e}(\mathbf{M};\mathbf{N}^{\prime \prime})\) ; if \(\omega\) and \(\omega^{\prime}\) are respectively sections of \(\mathsf{Alt}^{d}(\mathbf{M};\mathbf{N})\) and \(\mathsf{Alt}^{e}(\mathbf{M};\mathbf{N}^{\prime})\) over an open set U, the section \(u(\omega ,\omega^{\prime})\) of \(\mathsf{Alt}^{d + e}(\mathbf{M};\mathbf{N}^{\prime \prime})\) over U will be denoted \(\omega \wedge_{\varphi}\omega^{\prime}\) , and called the exterior product of \(\omega\) and \(\omega^{\prime}\) .

We have the formula:

\[
(1) \quad \left(\omega \wedge_ {\varphi} \omega^ {\prime}\right) \left(s _ {1}, \dots , s _ {d + e}\right) = \sum_ {\sigma} \varepsilon_ {\sigma} \varphi \left(\omega \left(s _ {\sigma (1)}, \dots , s _ {\sigma (d)}\right), \omega^ {\prime} \left(s _ {\sigma (d + 1)}, \dots , s _ {\sigma (d + e)}\right)\right)
\]

where the \(s_{i}\) are sections of M over U, and where the summation is extended over the permutations \(\sigma\) of \(\{1,2,\ldots,d+e\}\) such that

\[
\sigma (1) <   \dots <   \sigma (d) \quad \text { and } \quad \sigma (d + 1) <   \dots <   \sigma (d + e).
\]

7.8.3. Let M be a vector bundle and A a bundle of algebras, with base B. Suppose that the fibres \(A_{b}\) of A are associative and commutative algebras, possessing an identity element, denoted \(e_{b}\) . For every open set U of B, we denote by \(\Omega^{d}(U)\) the \(\mathcal{C}^{r}(U)\)-module formed by the sections of the bundle \(\operatorname{Alt}^{d}(M;A)\) and \(\Omega^{*}(U)\) the direct sum of the \(\Omega^{d}(U)\) for \(d\geqslant0\) . The multiplications on each fibre define a pairing of \(A\times_{B}A\) into A, whence (7.8.2) a structure of graded algebra on \(\Omega^{*}(U)\) , which is associative and anticommutative. The subalgebra \(\Omega^{0}(U)\) is the algebra of sections of A. An element \(\omega\) of \(\Omega^{1}(\mathbf{U})\) is identified with a U-morphism of \(\mathbf{M}|\mathbf{U}\) into \(\mathbf{A}|\mathbf{U}\) (7.7.3): if \(s\in\mathcal{S}_{\mathbf{M}}^{r}(\mathbf{U})\) , we denote by \(\langle\omega,s\rangle\) the section \(\omega(s)\) of A (7.4.2). Let \(s_{j}\in\mathcal{S}_{\mathbf{M}}^{r}(\mathbf{U})\) and \(\omega_{j}\in\Omega^{1}(\mathbf{U})\) (for \(1\leqslant j\leqslant d\) ); we have:

\(^{1}\) The reader should take care not to confuse this use of the letter \(\omega\) with that defined on p. 10.

\[
\omega (s _ {1}, \dots , s _ {d}) = \det (\langle \omega_ {i}, s _ {j} \rangle) \quad \text {   for   } \omega = \omega_ {1} \wedge \dots \wedge \omega_ {d}.\tag{2}
\]

7.8.4. Let \(d \geqslant 1\) . There exists a pairing \(i\) of \(\mathbf{M} \times_{\mathbf{B}} \mathrm{Alt}^{d}(\mathbf{M}; \mathbf{A})\) to \(\mathrm{Alt}^{d-1}(\mathbf{M}; \mathbf{A})\) whose restriction to each fiber is given by the right inner product (cf. A, III, p. 156). If \(s\) is a section of \(\mathbf{M}\) on the open subset \(\mathbf{U}\) and if \(\omega \in \Omega^{d}(\mathbf{U})\) , one denotes by \(i(s)\omega\) the section \(i(s, \omega)\) of \(\mathrm{Alt}^{d-1}(\mathbf{M}; \mathbf{A})\) on \(\mathbf{U}\) ; we set \(i(s)\omega = 0\) for \(\omega\) in \(\Omega^{0}(\mathbf{U})\) .

We thus associate to every section \(s\) of \(\mathbf{M}\) on \(\mathbf{U}\) an endomorphism of the \(\mathcal{C}^r (\mathbf{U})\)-module \(\Omega^{*}(\mathbf{U})\) . We have the formulas:

\[
(3) \quad (i (s) \omega) (s _ {1}, \dots , s _ {d - 1}) = \omega (s, s _ {1}, \dots , s _ {d - 1}) \quad \text { for } \omega \in \Omega^ {d} (\mathbf {U}), d \geqslant 1
\]

\begin{enumerate}
\def\labelenumi{(\arabic{enumi})}
\setcounter{enumi}{3}
\tightlist
\item
  \(i(s)\circ i(s) = 0\)
\end{enumerate}

\[
i (s) \omega = \langle \omega , s \rangle \quad \text { for } \omega \in \Omega^ {1} (\mathbf {U}) \tag {5}
\]

\[
(6) \quad i (s). (\omega \wedge \omega^ {\prime}) = i (s) \omega \wedge \omega^ {\prime} + (- 1) ^ {d} \omega \wedge i (s) \omega^ {\prime} \quad \text { for } \omega \in \Omega^ {d} (\mathbf {U})
\]

\[
i (s) (\omega_ {1} \wedge \dots \wedge \omega_ {p}) = \sum_ {i = 1} ^ {p} (- 1) ^ {i + 1} \langle \omega_ {i}, s \rangle \omega_ {1} \wedge \dots \wedge \hat {\omega} _ {i} \wedge \dots \wedge \omega_ {d}.\tag{7}
\]

In the last formula, the \(\omega_{i}\) are in \(\Omega^{1}(\mathbf{U})\) and the sign indicates that the symbol it surmounts is to be omitted.

All the operations described above on sections are multilinear over the ring \(\mathcal{C}^r (\mathbf{U};\mathbf{K})\) .

7.8.5. Let L be a Banach algebra over K. The definitions and results of Sections 7.7 and 7.8 extend to the case of vector bundles over L: one defines analogously the bundles of L-multilinear or alternating L-multilinear maps.

